\documentclass[11pt,reqno]{amsart}

\usepackage[T1]{fontenc}
\usepackage[utf8]{inputenc}
\usepackage{lmodern}
\usepackage{amsmath,amssymb,amsthm,mathtools}
\usepackage{booktabs,longtable,array,tabularx}
\usepackage{microtype}
\usepackage{geometry}
\usepackage{hyperref}
\usepackage{xurl}
\usepackage{enumitem}
\geometry{margin=1in}
\hypersetup{
  colorlinks=true,
  linkcolor=blue,
  citecolor=blue,
  urlcolor=blue,
  pdftitle={A Certified Contact-Geometric Lower Bound for the Three-Dimensional Blaschke--Lebesgue Problem},
  pdfauthor={Bruce J. Swihart},
  pdfsubject={Version v1.0.0; computer-assisted preprint; three-dimensional Blaschke--Lebesgue problem},
  pdfkeywords={constant width, Blaschke--Lebesgue problem, convex geometry, computer-assisted proof, exact rational arithmetic}
}

\newtheorem{theorem}{Theorem}[section]
\newtheorem{lemma}[theorem]{Lemma}
\newtheorem{proposition}[theorem]{Proposition}
\newtheorem{corollary}[theorem]{Corollary}
\theoremstyle{remark}
\newtheorem{remark}[theorem]{Remark}

\newcommand{\R}{\mathbb R}
\newcommand{\Sph}{\mathbb S^2}
\newcommand{\Vol}{\operatorname{Vol}}
\newcommand{\conv}{\operatorname{conv}}
\newcommand{\diam}{\operatorname{diam}}
\newcommand{\Id}{\operatorname{Id}}
\newcommand{\dd}{\,d\sigma}
\newcommand{\eps}{\varepsilon}
\newcommand{\releaseversion}{v1.0.0}
\newcommand{\releasedate}{September 26, 2026}
\newcommand{\repositoryurl}{https://github.com/swihart/minvol-degree3-contact-bound}
\newcommand{\releaseurl}{https://github.com/swihart/minvol-degree3-contact-bound/releases/tag/v1.0.0}

% Generated by paper/reconcile_certificate.py.  Do not edit by hand.
\providecommand{\CertificateSchema}{minvol.universal\_contact\_bound.v1}
\providecommand{\CertificateSha}{e6b834ccc35241e94074c0b4ecd9045e6c4383cf4185150c78a0f24e021b66ea}
\providecommand{\MainNumerator}{26221074253}
\providecommand{\MainDenominator}{200000000000}
\providecommand{\MainCoefficient}{\frac{26221074253\pi}{200000000000}}
\providecommand{\MainCoefficientOfPi}{\frac{26221074253}{200000000000}}
\providecommand{\MainDecimal}{0.4118796712122863775432315385981549188973}
\providecommand{\GainDecimal}{0.0000021074319841032950526126153150478947}
\providecommand{\SplitRadius}{\frac{305695601}{500000000}}
\providecommand{\SplitDecimal}{0.611391202}
\providecommand{\DeficiencyDecimal}{0.01939912533733453686615863855265}
\providecommand{\DeterminantDecimal}{0.99011227429940592416253966614320}
\providecommand{\FanCoefficientDecimal}{0.5851824745606993360954839072}
\providecommand{\CoreDecimal}{0.4117357103127771735489695298}
\providecommand{\CapsDecimal}{0.0001439609020552581838284512}
\providecommand{\NearDecimal}{0.4118796712148324317327979810}
\providecommand{\HandoffMarginDecimal}{0.0000000000025460541895664424}
\providecommand{\CrossAxisGapDecimal}{0.6378776497509377628400649501}
\providecommand{\WidthZeroDecimal}{0.995185216461759158}
\providecommand{\WidthOneDecimal}{0.998387315752931109}
\providecommand{\WidthTwoDecimal}{0.998387315752931109}
\providecommand{\CapZeroDecimal}{0.000050275601362298174}
\providecommand{\CapOneDecimal}{0.000046842650346480004}
\providecommand{\CapTwoDecimal}{0.000046842650346480004}
\providecommand{\HullZeroPlus}{110}
\providecommand{\HullZeroMinus}{84}
\providecommand{\HullOnePlus}{102}
\providecommand{\HullOneMinus}{76}
\providecommand{\HullTwoPlus}{102}
\providecommand{\HullTwoMinus}{76}
\providecommand{\FanVertices}{98{,}306}
\providecommand{\FanFaces}{196{,}608}
\providecommand{\FanEdges}{294{,}912}
\providecommand{\FanOrbits}{24{,}833}
\providecommand{\ChangedVertices}{18{,}229}
\providecommand{\FanMarginDecimal}{0.0000000002888904207092857876733697822441}


\title[Certified contact-geometric lower bound]{A Certified Contact-Geometric Lower Bound for the Three-Dimensional Blaschke--Lebesgue Problem}
\author{Bruce J. Swihart}
\thanks{No institutional affiliation is asserted. Correspondence and verification reports: \url{https://github.com/swihart/minvol-degree3-contact-bound/issues}.}
\date{Version \releaseversion\ -- \releasedate}
\subjclass[2020]{52A20, 52A40, 65G30}
\keywords{Bodies of constant width, Blaschke--Lebesgue problem, convex geometry, computer-assisted proof, exact rational arithmetic}

\begin{document}

\begin{abstract}
Let $K\subset\R^3$ be a convex body of constant width $d>0$.  We give an
exact-rational computer-assisted certificate for the universal estimate
\[
  \Vol(K)\ge \MainCoefficient\,d^3
  =\MainDecimal\ldots\,d^3.
\]
The proof divides the possible circumradii at the rational value
$R_*=\SplitRadius$.  Below $R_*$, a finite radial-propagation certificate
refines the contact-geometric method of HYRA.  Above $R_*$, four balanced
circumsphere contacts force a high-resolution polyhedral core; a sharp affine
determinant transfer, a dominant-edge coordinate-width estimate, and six
pairwise disjoint cone caps complete the argument.  Every claim-bearing
finite comparison is checked with integers and rational arithmetic.  A
binary64 Python implementation and a base-R implementation are independent
numerical audits, and an adversarial suite verifies rejection of seven classes
of corrupted proof objects.  This preprint does not address equality, sharpness,
the identity of a minimizer, or the Meissner conjecture.
\end{abstract}

\maketitle

\begin{center}
\fbox{\begin{minipage}{0.92\textwidth}
\small\textbf{Verification status.}
This is version \releaseversion\ of an AI-assisted,
computer-assisted research preprint.  The standalone exact verifier,
independent Python and base-R audits, reconstruction tests, and adversarial
mutation tests have passed in the accompanying package.  The canonical
repository is \url{\repositoryurl}, and the versioned release is
\url{\releaseurl}.  External mathematical reproduction and peer review are
not asserted.  The exact certificate JSON and exact verifier are the proof
authority for the finite computations.
\end{minipage}}
\end{center}

\section{Introduction}

A convex body has constant width $d$ when every pair of parallel supporting
planes orthogonal to a unit direction is separated by the same distance $d$.
In the plane, the Blaschke--Lebesgue theorem identifies the Reuleaux triangle
as the least-area body of prescribed constant width
\cite{Lebesgue1914,Blaschke1915}.  The analogous three-dimensional problem is
open.  The two classical Meissner bodies have common volume
\[
 V_M(d)=\frac{\pi}{12}\left(8-3\sqrt3\arccos\frac13\right)d^3
       =0.419860045965080\ldots d^3,
\]
and are conjectured to minimize volume; this explicit comparison value is not
a proved universal minimum \cite{KawohlWeber2011,Martini2019}.

Chakerian proved the lower bound
$\frac{\pi}{3}(3\sqrt6-7)d^3=0.364916122595\ldots d^3$
\cite{Chakerian1966}.  Nishioka later obtained
$\frac{4\pi}{33}d^3=0.380799109526\ldots d^3$ from a support-function
formulation, spherical harmonics, and Bochner's formula
\cite{Nishioka2026}.  A subsequent contact-geometric manuscript by HYRA
combined the minimum circumsphere, radial propagation, a divergence estimate,
and exact cap recovery to reach a certified coefficient above
$0.411040473721188$ \cite{HYRA2026}.  The present certificate develops that
contact-geometric route further.

The main claim of the package is the following.

\begin{theorem}[Main computer-assisted claim]\label{thm:main}
Every convex body $K\subset\R^3$ of constant width $d>0$ satisfies
\[
 \boxed{\Vol(K)\ge \MainCoefficient\,d^3
 =\MainDecimal\ldots\,d^3.}
\]
\end{theorem}

The argument is homogeneous, so we work at width one.  If $R$ is the radius of
a minimum circumscribed ball, Jung's theorem gives
\[
 \frac12\le R\le\frac{\sqrt6}{4}.
\]
The certificate uses the exact handoff
\[
 R_*=\SplitRadius=\SplitDecimal.
\]
A finite radial certificate covers $R\le R_*$.  The complementary interval is
handled by a nearly tetrahedral four-contact configuration, a
$\FanVertices$-ray forced fan, and six additional cone caps.  The near-Jung
branch gives the strictly larger lower value
$\NearDecimal\ldots$, leaving the final row of the radial partition as the
unique bottleneck.

The proof has two layers.  The geometric reductions are stated and proved in
the text.  The finite claims are represented by immutable integer and rational
objects and replayed by the exact verifier.  Section~\ref{sec:certificate}
records the trust boundary and the division between proof authority and
independent audits.

\section{Constant-width identities and minimum circumsphere contacts}
\label{sec:prelim}

Let $h_K$ be the support function of a width-one body $K$.  Then
\begin{equation}\label{eq:width-support}
 h_K(u)+h_K(-u)=1\qquad(u\in\Sph).
\end{equation}
Consequently
\begin{equation}\label{eq:difference-body}
 K+(-K)=B(0,1),\qquad \diam K=1.
\end{equation}
In dimension three, Blaschke's identity relates volume and surface area:
\begin{equation}\label{eq:blaschke}
 S(K)=2\Vol(K)+\frac{2\pi}{3}.
\end{equation}
These standard identities are reviewed in
\cite{ChakerianGroemer1983,Martini2019,Schneider2014}.

Translate a minimum circumscribed ball to $B(0,R)$ and let
$F=K\cap\partial B(0,R)$.  Minimality of the ball implies
$0\in\conv F$.  Caratheodory's theorem therefore gives contact directions
$u_i\in\Sph$ and positive weights $\lambda_i$ such that
\begin{equation}\label{eq:contact-balance}
 p_i=Ru_i\in K,\qquad
 \sum_i\lambda_i=1,\qquad
 \sum_i\lambda_i p_i=0,
\end{equation}
using at most four contacts.  Since $\diam K=1$,
$|p_i-p_j|\le1$.  The usual balance argument yields $R^2\le3/8$.  When
$R^2>1/3$, four contacts are necessary; this is the regime used in the
near-Jung branch.

Constant width also supplies an opposite point at each support point.

\begin{lemma}[Opposite-point identity]\label{lem:opposite}
If $x\in\partial K$ has an outer unit normal $n$, then $x-n\in K$.  Hence,
when $K\subset B(0,R)$,
\begin{equation}\label{eq:opposite-ineq}
 \langle x,n\rangle\ge \frac{|x|^2+1-R^2}{2}.
\end{equation}
\end{lemma}

\begin{proof}
Let $y$ lie on the supporting plane opposite $x$ in direction $n$.  Width one
gives $\langle x-y,n\rangle=1$, whereas \eqref{eq:difference-body} gives
$|x-y|\le1$.  Equality in Cauchy--Schwarz forces $x-y=n$.  Since
$x-n\in K\subset B(0,R)$, expanding $|x-n|^2\le R^2$ gives
\eqref{eq:opposite-ineq}.
\end{proof}

Because $R<1$, the circumcenter lies in the interior of $K$.  We write
$\rho:\Sph\to(0,\infty)$ for the radial function about that center.

\section{The lower-circumradius radial certificate}
\label{sec:lower}

This section records the analytic reduction and the finite exact data used on
$[1/2,R_*]$.  The construction follows the contact-geometric framework in
\cite{HYRA2026}, with a refined terminal partition.

\subsection{Transport inequality and tangent gap}

For $0<m\le1-R^2$, define
\[
 F_m(x)=\frac{2x}{|x|^2+m}.
\]
By Lemma~\ref{lem:opposite}, $\langle F_m(x),n\rangle\ge1$ on every smooth
support point.  Applying the divergence theorem and then spherical
integration gives
\begin{equation}\label{eq:transport}
 3V=\int_{\Sph}\rho^3\dd,
 \qquad
 S\le\int_{\Sph}\frac{2\rho^3}{\rho^2+m}\dd,
\end{equation}
first for $C^1$ bodies and then for arbitrary constant-width bodies by outer
parallel regularization.

Put
\[
 \varphi_m(X)=\frac{2X}{X^{2/3}+m}\qquad(X>0).
\]
A direct differentiation shows $\varphi_m''<0$.  If $aX+b$ is the tangent at
$X=s_0^3$, the gap
\[
 G_{m,s_0}(s)=as^3+b-\frac{2s^3}{s^2+m}
\]
is nonnegative, vanishes only at $s_0$, and is monotone away from $s_0$.
Writing $E_G=\int_{\Sph}G_{m,s_0}(\rho)\dd$, equations
\eqref{eq:blaschke} and \eqref{eq:transport} imply
\begin{equation}\label{eq:master}
 V\ge\frac{2\pi/3-4\pi b+E_G}{3a-2},
 \qquad 3a-2>0.
\end{equation}
Thus the purely analytic bound discards $E_G\ge0$, while a stronger finite
certificate recovers part of the gap on forced spherical caps.

\subsection{Rational radial propagation}

Suppose $Ru_0\in K$.  Diameter one implies the antipodal ball estimate
\begin{equation}\label{eq:antipodal-ball}
 \rho(v)\le -R_*c+\sqrt{1-R_*^2(1-c^2)},
 \qquad c=-\langle v,u_0\rangle\in[0,1],
\end{equation}
whenever $R_*\le R$.  In the other direction, the supporting-plane bound
\eqref{eq:opposite-ineq} propagates a radial lower estimate from one direction
to a nearby direction.  With
\[
 \cos\gamma_m(s)=\frac{s^2+m}{2s},
\]
the one-step upper propagation map is
\begin{equation}\label{eq:propagation-map}
 N_m(s,\delta)=
 \frac{s(s^2+m)}{(s^2+m)\cos\delta-
 \sqrt{4s^2-(s^2+m)^2}\sin\delta}.
\end{equation}
The certificate uses the contrapositive of this estimate to construct radial
lower chains away from each contact.

Every angle is represented by a rational tangent-half-angle grid:
\begin{equation}\label{eq:rational-angle}
 t_k=\frac{k}{N},\qquad
 \cos\theta_k=\frac{1-t_k^2}{1+t_k^2},\qquad
 \sin\theta_k=\frac{2t_k}{1+t_k^2}.
\end{equation}
After clearing the one square root in \eqref{eq:propagation-map}, each
propagation step is reduced to signed integer polynomial inequalities.  The
same is true for \eqref{eq:antipodal-ball}, cap disjointness, and the tangent-gap
sum.  The resulting arrays therefore contain no floating-point proof step.

\subsection{Finite partition}

The exact lower branch contains twelve rows.  Four doubled-angular-grid cells
replace a formerly insufficient parent row; their angular denominator is
$76{,}800$ and their radial denominator is $4\cdot10^{10}$.  The weakest
replacement margin over the terminal coefficient of $\pi$ is
$1.8799274348\times10^{-5}$.  Appendix~\ref{app:lower-table} gives the complete
gap-free table generated directly from the certificate JSON.  Its unique
weakest row is \texttt{O3B02}, ending exactly at $R_*$.

\section{Near-Jung contact geometry}
\label{sec:contact}

Assume now $R\ge R_*$ and use four contacts in
\eqref{eq:contact-balance}.  Put $q=R^2$ and define the six edge deficiencies
\begin{equation}\label{eq:deficiency}
 e_{ij}=1-|p_i-p_j|^2\ge0,
 \qquad S=\sum_{i<j}e_{ij}.
\end{equation}
The Gram-kernel relation is equivalent to
\begin{equation}\label{eq:weighted-balance}
 \lambda_i+\sum_{j\ne i}\lambda_j e_{ij}=1-2q.
\end{equation}
Summing over $i$ gives
\begin{equation}\label{eq:weighted-sum}
 \sum_{i<j}(\lambda_i+\lambda_j)e_{ij}=3-8q.
\end{equation}
Equation \eqref{eq:weighted-balance} implies $\lambda_i\le1-2q$; hence
$\lambda_i+\lambda_j\ge4q-1$.  Therefore
\begin{equation}\label{eq:scalar-budget}
 \boxed{S\le E(R):=\frac{3-8R^2}{4R^2-1}.}
\end{equation}
The function $E(R)$ decreases on the Jung interval, so throughout the
near-Jung branch
\[
 S\le E_*:=E(R_*)=\DeficiencyDecimal\ldots.
\]

\subsection{Affine reference coordinates}

Let
\[
 v_0=(1,1,1),\quad v_1=(1,-1,-1),\quad
 v_2=(-1,1,-1),\quad v_3=(-1,-1,1).
\]
There are an invertible linear map $L$ and a balance-center vector $c$ such
that
\[
 p_i=L(v_i-c).
\]
Set $A=2\sqrt2L$ and $H=A^TA=8L^TL$.  The regular unit-edge tetrahedron
corresponds to $H=I$.  Solving the six edge equations gives
\begin{align}
 H_{00}&=1+\tfrac12(e_{01}-e_{02}-e_{03}-e_{12}-e_{13}+e_{23}),\label{eq:h00}\\
 H_{11}&=1+\tfrac12(-e_{01}+e_{02}-e_{03}-e_{12}+e_{13}-e_{23}),\label{eq:h11}\\
 H_{22}&=1+\tfrac12(-e_{01}-e_{02}+e_{03}+e_{12}-e_{13}-e_{23}).\label{eq:h22}
\end{align}

The affine volume loss is controlled by a sharp scalar determinant estimate.

\begin{proposition}[Determinant transfer]\label{prop:determinant}
For nonnegative deficiencies with $S<3/2$,
\begin{equation}\label{eq:determinant-bound}
 \det H\ge1-\frac S2-\frac{S^2}{2}.
\end{equation}
Consequently, on the near-Jung branch,
\[
 \det H\ge\DeterminantDecimal\ldots.
\]
\end{proposition}

\begin{proof}
Use three edge vectors from one contact and write their normalized Gram matrix
in terms of the six squared edge lengths $1-e_{ij}$.  Twice its determinant is
$\det H$.  Expanding exactly and subtracting the right side of
\eqref{eq:determinant-bound} produces fifteen quadratic cross terms, each with
coefficient at least one, together with twelve negative squarefree cubic terms,
each with coefficient $-1/2$, and additional nonnegative cubic terms.  Each
pair of deficiency variables occurs in at most four of the negative triples.

For a negative triple $abc$,
\[
 abc\le\frac{S}{3}(ab+ac+bc).
\]
The total negative cubic part is therefore at least
$-(2S/3)\sum_{r<s}e_re_s$, whereas the positive quadratic part is at least
$\sum_{r<s}e_re_s$.  Since $S<3/2$, the difference is nonnegative.  The exact
verifier reconstructs the polynomial and checks the complete coefficient and
pair-incidence census rather than trusting a stored expansion.
\end{proof}

\section{The high-resolution forced fan}
\label{sec:fan}

The near-Jung core is represented in the regular reference coordinates.  The
unit sphere is subdivided by the cubical grid of resolution $N=128$.  After
identifying shared face boundaries, the triangulation has
\[
 \FanVertices\text{ rays},\qquad
 \FanFaces\text{ triangular cones},\qquad
 \FanEdges\text{ edges}.
\]
The order-four stabilizer of the distinguished contact edge has
$\FanOrbits$ ray orbits.

For every ray, the certificate stores an integer radial numerator and a dyadic
S-procedure multiplier.  The exact verifier reconstructs the grid, its
triangulation, and the stabilizer orbits from first principles.  It then checks
every ray inequality over the full nonnegative deficiency simplex
\[
 e_{01}\ge e_{ij},\qquad \sum e_{ij}\le E_*.
\]
The raywise optimization is exact: its candidates are the vertices forced by
the dominant-edge budget, and all remaining arithmetic is rational.  Relative
to the immutable reference fan, $\ChangedVertices$ radial numerators move
inward while all inherited multiplier numerators remain fixed.  The least
nontrivial exact margin is
\[
 \FanMarginDecimal\ldots>10^{-12}.
\]

Let $Q_i\in\mathbb Z^3$ be the reconstructed cubical rays and let $r_i$ be the
certified rational radii.  Summing the signed determinants
$\det(Q_i,Q_j,Q_k)r_ir_jr_k/6$ over the fixed cone triangulation gives the
reference coefficient
\[
 C_{\rm fan}=\FanCoefficientDecimal\ldots.
\]
The affine determinant factor from Proposition~\ref{prop:determinant} then
gives the physical forced-core estimate
\begin{equation}\label{eq:core-volume}
 \boxed{V_{\rm core}\ge\CoreDecimal\ldots.}
\end{equation}
The stored volume numerator is independently reconstructed from all
$\FanFaces$ cones.

\section{Dominant-edge coordinate widths}
\label{sec:width}

Relabel the contacts so that $e_{01}$ is a largest deficiency.  The fan atlas
includes ties, so this does not discard any configuration.

\begin{proposition}[Coordinate-specific transformed widths]\label{prop:widths}
If $e_{01}\ge e_{ij}$ for every edge, then
\[
 H_{00}\le1+\frac S2,
 \qquad
 H_{11},H_{22}\le1+\frac S6.
\]
If $g_j=\sqrt{(H^{-1})_{jj}}$, then
\begin{equation}\label{eq:width-lower}
 g_0\ge\frac1{\sqrt{1+S/2}},
 \qquad
 g_1,g_2\ge\frac1{\sqrt{1+S/6}}.
\end{equation}
\end{proposition}

\begin{proof}
The first diagonal estimate follows immediately from \eqref{eq:h00}.  For the
second,
\[
 S-3[2(H_{11}-1)]
 =2(e_{01}-e_{02})+2(e_{01}-e_{13})
  +4(e_{03}+e_{12}+e_{23})\ge0.
\]
The identity for $H_{22}$ is obtained by interchanging the corresponding
indices.  Proposition~\ref{prop:determinant} makes $H$ positive definite, and
Cauchy--Schwarz in the $H$-inner product gives
$1\le H_{jj}(H^{-1})_{jj}$.
\end{proof}

Using $S\le E_*$ and exact rational lower enclosures for the square roots, the
certificate obtains
\[
 (g_0,g_1,g_2)\ge
 (\WidthZeroDecimal\ldots,\WidthOneDecimal\ldots,\WidthTwoDecimal\ldots).
\]
The distinction between the dominant direction and the two transverse
directions is essential for the cap contribution.

\section{Six optimized cone caps}
\label{sec:caps}

For coordinate axis $j$, let $h_{j,+}$ and $h_{j,-}$ be the exact signed
supports of the reference fan.  Intersecting the fan with fixed rational planes
below these supports gives exact convex section polygons of areas
$A_{j,+}$ and $A_{j,-}$.  If $t_{j,+}$ and $t_{j,-}$ are the remaining exterior
support extensions, Proposition~\ref{prop:widths} gives
\begin{equation}\label{eq:total-extension}
 t_{j,+}+t_{j,-}\ge
 T_j:=g_j-\frac{h_{j,+}+h_{j,-}}{2\sqrt2}.
\end{equation}

A section of physical area $A$, lying a depth $d$ below the fan support, and
an exterior support extension $t$ force a cone of volume at least
\begin{equation}\label{eq:cap-profile}
 \phi_{A,d}(t)=\frac{At^3}{24(d+t)^2}.
\end{equation}
For $t>0$,
\[
 \phi_{A,d}''(t)=\frac{Ad^2t}{4(d+t)^4}>0.
\]
Thus, subject to equality in \eqref{eq:total-extension}, each antipodal pair
has a unique minimizing allocation.  The exact verifier brackets each
stationary point by rationals at distance at most $10^{-20}$ and proves
opposite derivative signs at the two bracket endpoints.

\begin{table}[ht]
\centering
\caption{Certified physical cone-cap contributions at the handoff radius.}
\label{tab:caps}
\begin{tabular}{@{}c c c c@{}}
\toprule
Axis & Width lower & Section vertices $+/-$ & Pair contribution\\
\midrule
$0$ & $\WidthZeroDecimal\ldots$ & $\HullZeroPlus/\HullZeroMinus$ & $\CapZeroDecimal\ldots$\\
$1$ & $\WidthOneDecimal\ldots$ & $\HullOnePlus/\HullOneMinus$ & $\CapOneDecimal\ldots$\\
$2$ & $\WidthTwoDecimal\ldots$ & $\HullTwoPlus/\HullTwoMinus$ & $\CapTwoDecimal\ldots$\\
\midrule
\multicolumn{3}{r}{Total} & $\CapsDecimal\ldots$\\
\bottomrule
\end{tabular}
\end{table}

The two caps associated with one axis lie in opposite support halfspaces.  For
distinct axes, the balance-center uncertainty and the matrix inequality
$H\succeq(1-E_*)I$ show that an overlap would force two signed coordinates to
exceed the least fan support while every transformed body point remains in a
fixed ellipsoid.  The exact contradiction gap is
\[
 \CrossAxisGapDecimal\ldots>0.
\]
Hence all six cap interiors are pairwise disjoint; possible boundary contacts
have zero volume.

\section{Universal assembly}
\label{sec:assembly}

Adding \eqref{eq:core-volume} and the three disjoint antipodal cap-pair
contributions gives
\[
 V_{\rm near}\ge\NearDecimal\ldots.
\]
The exact margin over the upper rational enclosure of the coefficient in
Theorem~\ref{thm:main} is
\[
 \HandoffMarginDecimal\ldots>0.
\]
Therefore the near-Jung construction covers
$[R_*,\sqrt6/4]$.

For $R\le R_*$, the exact partition in Section~\ref{sec:lower} applies.  Its
unique weakest row is \texttt{O3B02}, whose exact coefficient is
$\MainCoefficientOfPi\pi$.  The two branch endpoints agree exactly, so the
union is the full Jung interval.  Finally, dilation from width one to width
$d$ multiplies volume by $d^3$, completing the computer-assisted proof of
Theorem~\ref{thm:main} under the trust boundary below.

\section{Certificate architecture and trust boundary}
\label{sec:certificate}

The public-format proof package is standalone: no path points into the private
research repository, and the complete replay begins from the files committed
with this manuscript.

\paragraph{Exact proof authority.}
The program \texttt{certificate/certify\_universal\_bound.py} reconstructs the
claim using Python integers and \texttt{Fraction} arithmetic.  NumPy is used to
read immutable integer arrays, not to decide theorem inequalities.  The
verifier checks the lower radial rows, reconstructs the fan and all cone
volumes, verifies the determinant census and coordinate-width identities,
rebuilds the six section polygons, proves the cap-allocation brackets and
six-cap disjointness, checks the branch handoff, and regenerates the certificate
JSON byte for byte.

\paragraph{Independent audits.}
A separately organized NumPy/binary64 program reconstructs the geometry
numerically.  A base-R program checks the archived finite decimals and
structural identities without importing Python.  Neither audit is used as proof
authority.  The Python test suite also reconstructs theorem-critical objects,
and seven adversarial tests alter one lower-row endpoint, one fan radius, one
multiplier, one section-hull vertex, one cap bracket, one disjointness threshold,
and one theorem digit.  Every altered object is rejected.

\paragraph{Paper reconciliation.}
The displayed certificate-critical constants and Appendix~\ref{app:lower-table}
are generated from the machine-readable JSON by
\texttt{paper/reconcile\_certificate.py}.  A stale or manually divergent TeX
include fails verification.  The certificate used for this manuscript has
SHA-256
\[
 \texttt{\CertificateSha}.
\]

\paragraph{Trusted components.}
The finite proof trusts the semantics of CPython integer arithmetic,
\path{fractions.Fraction}, JSON and NPZ parsing, the small exact verifier
source, and the operating system's faithful reading of checksum-protected
files.  It does not trust binary floating point for a theorem inequality.
Typesetting software and generated PDF bytes are outside the mathematical trust
boundary.

\begin{table}[ht]
\centering
\caption{Principal certificate census.}
\label{tab:census}
\begin{tabular}{@{}lr@{}}
\toprule
Object & Count or certified value\\
\midrule
Circumradius split & $\SplitRadius$\\
Lower-branch rows & $12$\\
Fan rays & $\FanVertices$\\
Fan triangular cones & $\FanFaces$\\
Fan stabilizer orbits & $\FanOrbits$\\
Changed fan vertices & $\ChangedVertices$\\
Minimum exact fan margin & $\FanMarginDecimal\ldots$\\
Exact section hulls & $6$\\
Antipodal allocation brackets & $3$\\
Near-Jung handoff margin & $\HandoffMarginDecimal\ldots$\\
\bottomrule
\end{tabular}
\end{table}

\section{Scope and limitations}
\label{sec:limitations}

The result is a universal lower-bound certificate.  It does not prove that the
coefficient is sharp, characterize equality, identify a minimizing body, or
establish Meissner extremality.  The explicit Meissner volume remains an upper
bound for the unknown minimum, not a consequence of this certificate.  The
paper also makes no claim of peer review or external independent reproduction.

The present result combines a human-readable geometric reduction with a finite
exact computation.  Future work may seek stronger circumradius-local estimates
or alternative certificates, but no such refinement is used in the theorem
proved here.

\section*{Acknowledgments and disclosure}

This manuscript and its verification package were developed with extensive
assistance from OpenAI ChatGPT (GPT-5.6 Sol Pro) during August--September 2026.
The AI system is not an author, reviewer, or proof authority.  Bruce J. Swihart
is the sole named author and is responsible for the claims, source selection,
repository contents, release decisions, and correction of errors.  No
institutional affiliation is asserted.  Detailed provenance is recorded in
\path{AI_ASSISTANCE.md}.  The paper and prose are licensed CC BY 4.0; software
and build infrastructure are licensed MIT.  The intended canonical repository
is \url{\repositoryurl}.

\begin{thebibliography}{99}

\bibitem{Blaschke1915}
W.~Blaschke,
\emph{Konvexe Bereiche gegebener konstanter Breite und kleinsten Inhalts},
Math. Ann. 76 (1915), 504--513.

\bibitem{Chakerian1966}
G.~D. Chakerian,
\emph{Sets of constant width},
Pacific J. Math. 19 (1966), no.~1, 13--21.

\bibitem{ChakerianGroemer1983}
G.~D. Chakerian and H.~Groemer,
\emph{Convex bodies of constant width},
in \emph{Convexity and its Applications}, Birkh\"auser, 1983, 49--96.

\bibitem{Harrell2002}
E.~M. Harrell II,
\emph{A direct proof of a theorem of Blaschke and Lebesgue},
J. Geom. Anal. 12 (2002), no.~1, 81--88.

\bibitem{HYRA2026}
HYRA,
\emph{A certified geometric lower bound for the three-dimensional
Blaschke--Lebesgue problem}, manuscript dated August 15, 2026.

\bibitem{Jung1901}
H.~Jung,
\emph{Ueber die kleinste Kugel, die eine r\"aumliche Figur einschliesst},
J. Reine Angew. Math. 123 (1901), 241--257.

\bibitem{KawohlWeber2011}
B.~Kawohl and C.~Weber,
\emph{Meissner's mysterious bodies},
Math. Intelligencer 33 (2011), no.~3, 94--101.

\bibitem{LachandOudet2007}
T.~Lachand--Robert and E.~Oudet,
\emph{Bodies of constant width in arbitrary dimension},
Math. Nachr. 280 (2007), no.~7, 740--750.

\bibitem{Lebesgue1914}
H.~Lebesgue,
\emph{Sur le probl\`eme des isop\'erim\`etres et sur les domaines de largeur constante},
Bull. Soc. Math. France 42 (1914), 72--76.

\bibitem{Martini2019}
H.~Martini, L.~Montejano, and D.~Oliveros,
\emph{Bodies of Constant Width: An Introduction to Convex Geometry with
Applications}, Birkh\"auser/Springer, 2019.

\bibitem{Nishioka2026}
A.~Nishioka,
\emph{An improved lower bound for the three-dimensional
Blaschke--Lebesgue problem from spectral and dual perspectives},
J. Math. Anal. Appl. 566 (2027), no.~1, article 131038;
\href{https://doi.org/10.1016/j.jmaa.2026.131038}{doi:10.1016/j.jmaa.2026.131038};
\href{https://arxiv.org/abs/2606.01754}{arXiv:2606.01754}.

\bibitem{Schneider2014}
R.~Schneider,
\emph{Convex Bodies: The Brunn--Minkowski Theory}, 2nd expanded ed.,
Cambridge University Press, 2014.

\end{thebibliography}

\appendix

\section{Complete lower-circumradius partition}
\label{app:lower-table}

% Generated by paper/reconcile_certificate.py.  Do not edit by hand.
\begingroup
\renewcommand{\arraystretch}{1.55}
\setlength{\extrarowheight}{1.5pt}
\begin{longtable}{@{}lll r@{}}
\caption{Gap-free lower-circumradius certificate.  The final column is the displayed width-normalized volume coefficient; exact coefficients of $\pi$ are stored in the certificate JSON.}\label{tab:lower-partition}\\
\toprule
Case & Circumradius interval & Exact mechanism & Volume lower bound\\
\midrule
\endfirsthead
\toprule
Case & Circumradius interval & Exact mechanism & Volume lower bound\\
\midrule
\endhead
\midrule
\multicolumn{4}{r}{Continued on next page}\\
\endfoot
\bottomrule
\endlastfoot
\texttt{0} & $[\frac{1}{2},\frac{3}{5}]$ & analytic tangent bound & $0.417357494296972$\\
\texttt{O1A} & $[\frac{3}{5},\frac{6047}{10000}]$ & exact radial caps & $0.416797024215342$\\
\texttt{O1B} & $[\frac{6047}{10000},\frac{3047}{5000}]$ & exact radial caps & $0.412421390558881$\\
\texttt{O2A} & $[\frac{3047}{5000},\frac{61019}{100000}]$ & exact radial caps & $0.412690926230586$\\
\texttt{O2B} & $[\frac{61019}{100000},\frac{30549}{50000}]$ & exact radial caps & $0.411993918231919$\\
\texttt{O3A1R1} & $[\frac{30549}{50000},\frac{30553}{50000}]$ & doubled angular grid & $0.412144899831404$\\
\texttt{O3A1R2} & $[\frac{30553}{50000},\frac{30557}{50000}]$ & doubled angular grid & $0.412075419676059$\\
\texttt{O3A1R3} & $[\frac{30557}{50000},\frac{30561}{50000}]$ & doubled angular grid & $0.412005994020500$\\
\texttt{O3A1R4} & $[\frac{30561}{50000},\frac{6113}{10000}]$ & doubled angular grid & $0.411936623442487$\\
\texttt{O3A2} & $[\frac{6113}{10000},\frac{244541}{400000}]$ & exact radial caps & $0.411898771358480$\\
\texttt{O3B01} & $[\frac{244541}{400000},\frac{76421809}{125000000}]$ & exact radial caps & $0.411890321028654$\\
\textbf{\texttt{O3B02}} & $[\frac{76421809}{125000000},\frac{305695601}{500000000}]$ & exact radial caps & $\textbf{0.411877563780302}$\\
\end{longtable}
\endgroup


The exact row coefficients, angular and radial grids, positive and negative
step counts, tangent radii, and the four refined proof arrays are stored under
\texttt{certificate/certificates/}.  Decimal entries in
Table~\ref{tab:lower-partition} are generated displays; the verifier compares
exact rational coefficients of $\pi$.

\section{Paper-facing exact constants}
\label{app:constants}

For convenient review, the principal width-normalized quantities are
\begin{align*}
R_*&=\SplitRadius=\SplitDecimal,\\
E_*&=\DeficiencyDecimal\ldots,\\
\det H&\ge\DeterminantDecimal\ldots,\\
C_{\rm fan}&=\FanCoefficientDecimal\ldots,\\
V_{\rm core}&\ge\CoreDecimal\ldots,\\
V_{\rm caps}&\ge\CapsDecimal\ldots,\\
V_{\rm near}&\ge\NearDecimal\ldots,\\
V_{\rm universal}&=\MainCoefficient
=\MainDecimal\ldots.
\end{align*}
Every displayed value in this list is regenerated from the certificate JSON.

\section{Machine-readable proof objects}
\label{app:objects}

\begingroup
\footnotesize
\renewcommand{\arraystretch}{1.08}
\begin{tabularx}{\dimexpr\textwidth-14pt\relax}{@{}>{\raggedright\arraybackslash}p{0.43\textwidth}X@{}}
\toprule
Path & Role\\
\midrule
\path{universal_contact_bound_certificate.json}
& Theorem coefficient, exact branch data, fan and cap summaries, and claim boundary.\\
\path{refined_lower_cell_1.npz}--\path{refined_lower_cell_4.npz}
& Exact rational propagation chains replacing the insufficient lower-row cell.\\
\path{universal_contact_fan.npz}
& Integer radial data and dyadic multiplier numerators for the adjusted fan.\\
\path{reference_near_jung_fan.npz}
& Regression fan used to check nesting and inherited multipliers.\\
\path{all_axis_pair_section_hulls.tsv}
& Six exact section polygons used in the cone-cap calculation.\\
\path{fan_orbit_margin_audit.tsv}
& One exact margin record for each stabilizer orbit.\\
\path{audit_constants.csv}
& Exact/decimal bridge for the independent numerical audits.\\
\bottomrule
\end{tabularx}
\endgroup

The repository checksum manifests authenticate these objects before and after
replay.  The exact verifier also rebuilds derived TSV, CSV, and JSON objects
byte for byte.

\end{document}
